\subsubsection{Matrix Classes and Spectral Invariants}

Matrix spectral statements are operator statements written in the standard coordinates of $\mathbb{F}^n$. The coordinate viewpoint is useful because determinants, traces, transposes, and block patterns become visible.

\begin{definition} \label{def:matrix-characteristic-polynomial} \hlt{Characteristic Polynomial}\\
For $A\in\operatorname{Mat}(n,n,\mathbb{F})$, the \hlt{characteristic polynomial} of $A$ is $p_A(z)\defeq\det(zI_n-A)$.
\end{definition}

\begin{theorem} \label{thm:matrix-characteristic-polynomial-basic-coefficients} \hlt{Basic Form of the Characteristic Polynomial}\\
Let $A=[a_{ij}]\in\operatorname{Mat}(n,n,\mathbb{F})$. Then $p_A$ is monic of degree $n$ and $p_A(z)=z^n-(\operatorname{tr}A)z^{n-1}+\cdots+(-1)^n\det A$.\\
Moreover, for $\lambda\in\mathbb{F}$, $\lambda\in\operatorname{spec}(A)$ if and only if $p_A(\lambda)=0$.
\begin{proof}
The only determinant term of degree greater than $n-2$ comes from the identity permutation:
\begin{align}
(z-a_{11})(z-a_{22})\cdots(z-a_{nn})=z^n-(a_{11}+\cdots+a_{nn})z^{n-1}+\cdots.
\nonumber
\end{align}
Thus $p_A$ is monic and the coefficient of $z^{n-1}$ is $-\operatorname{tr}A$; also $p_A(0)=\det(-A)=(-1)^n\det A$. Finally, $p_A(\lambda)=0$ if and only if $\lambda I_n-A$ is singular, equivalently $\ker(A-\lambda I_n)\neq\{\mathbf{0}\}$.\\
\hlt{Structural Commentary}: The diagonal product determines the top two coefficients.\\
\hlt{Key Insight}: Eigenvalues are precisely the roots of $\det(zI_n-A)$.
\end{proof}
\end{theorem}

\begin{definition} \label{def:matrix-algebraic-multiplicity} \hlt{Algebraic Multiplicity}\\
For $\lambda\in\mathbb{F}$, let $m_A(\lambda)$ be the multiplicity of $\lambda$ as a root of $p_A$, with $m_A(\lambda)=0$ if $\lambda\notin\operatorname{spec}(A)$.\\
For $\lambda\in\operatorname{spec}(A)$, this is the \hlt{algebraic multiplicity} of $\lambda$. Unqualified multiplicity means algebraic multiplicity.
\end{definition}

\begin{figure}[H]
  \centering
  \begin{tikzpicture}[
  >=Stealth,
  node distance=0.85cm and 1.0cm,
  box/.style={draw=blue!55!black, fill=blue!5, rounded corners=2pt, align=center, minimum height=0.78cm, minimum width=2.05cm, font=\small},
  root/.style={circle, draw=red!70!black, fill=red!10, inner sep=1.5pt, font=\small},
  arr/.style={->, line width=0.65pt, draw=black!70}
]
  \node[box] (matrix) {$A$};
  \node[box, right=of matrix] (poly) {$p_A$};
  \node[box, right=of poly] (roots) {roots};

  \draw[arr] (matrix) -- (poly);
  \draw[arr] (poly) -- (roots);

  \node[right=0.72cm of roots, inner sep=0pt, font=\small] (openroots) {$\{$};
  \node[root, right=0.10cm of openroots] (l1) {$\lambda_1$};
  \node[root, right=0.24cm of l1] (l2) {$\lambda_2$};
  \node[root, right=0.24cm of l2] (ld) {$\lambda_d$};
  \node[right=0.10cm of ld, inner sep=0pt, font=\small] (closeroots) {$\}$};
  \draw[arr] (roots) -- (openroots);
\end{tikzpicture}

  \caption{Form $p_A(z)=\det(zI_n-A)$ and factor it; each root $\lambda$ is counted by $m_A(\lambda)$.}
  \label{fig:characteristic-polynomial-roots}
\end{figure}

\begin{proposition} \label{prop:upper-triangular-characteristic-polynomial} \hlt{Upper-Triangular Eigenvalues}\\
If $A=[a_{ij}]\in\operatorname{Mat}(n,n,\mathbb{F})$ is upper triangular, then $p_A(z)=\prod\limits_{j=1}^{n}(z-a_{jj})$.\\
Thus $\operatorname{spec}(A)=\{a_{jj}:1\leq j\leq n\}$, and for $\lambda\in\operatorname{spec}(A)$, $m_A(\lambda)=\#\{j:a_{jj}=\lambda\}$.
\begin{proof}
The matrix $zI_n-A$ is upper triangular with diagonal entries $z-a_{11},\ldots,z-a_{nn}$. Its determinant is the product of those diagonal entries. The roots of this product are precisely the diagonal entries, with repetition.\\
\hlt{Structural Commentary}: Triangular form removes all off-diagonal determinant contributions.\\
\hlt{Key Insight}: Once a matrix is triangular, the spectrum is visible on the diagonal.
\end{proof}
\end{proposition}

\begin{proposition} \label{prop:transpose-conjugate-adjoint-characteristic-polynomial} \hlt{Transpose, Conjugate, and Adjoint Spectra}\\
Let $A\in\operatorname{Mat}(n,n,\mathbb{C})$. Then $p_{A^t}(z)=p_A(z)$, $p_{\overline{A}}(z)=\overline{p_A(\overline{z})}$, and $p_{A^*}(z)=\overline{p_A(\overline{z})}$.\\
Consequently, if $\lambda\in\operatorname{spec}(A)$, then $m_{A^t}(\lambda)=m_A(\lambda)$ and $m_{\overline{A}}(\overline{\lambda})=m_{A^*}(\overline{\lambda})=m_A(\lambda)$.
\begin{proof}
Use determinant identities:
\begin{align}
p_{A^t}(z)&=\det(zI_n-A^t)=\det((zI_n-A)^t)=p_A(z), \nonumber\\
p_{\overline{A}}(z)&=\det(zI_n-\overline{A})
=\det(\overline{\overline{z}I_n-A})
=\overline{\det(\overline{z}I_n-A)}
=\overline{p_A(\overline{z})}.
\nonumber
\end{align}
Since $A^*=\overline{A}^{\,t}$, the first identity applied to $\overline{A}$ gives the formula for $A^*$.\\
Multiplicity is preserved because conjugating a polynomial sends a factor $(z-\lambda)^k$ to $(z-\overline{\lambda})^k$.\\
\hlt{Structural Commentary}: Transpose preserves $p_A$; conjugation conjugates its roots.\\
\hlt{Key Insight}: The adjoint combines conjugation with transposition, so its spectrum is conjugated.
\end{proof}
\end{proposition}

\begin{proposition} \label{prop:transpose-preserves-spectrum-and-geometric-multiplicity} \hlt{Transpose Preserves Spectrum and Geometric Multiplicity}\\
Let $A\in\operatorname{Mat}(n,n,\mathbb{F})$ and $\lambda\in\mathbb{F}$. Then $\lambda\in\operatorname{spec}(A)$ $\Longleftrightarrow$ $\lambda\in\operatorname{spec}(A^t)$, and $\dim E_{\lambda}(A)=\dim E_{\lambda}(A^t)$.
\begin{proof}
Since $(A-\lambda I_n)^t=A^t-\lambda I_n$, Theorem \ref{thm:row-rank-equals-column-rank} gives $\operatorname{rank}(A-\lambda I_n)=\operatorname{rank}(A^t-\lambda I_n)$.\\
Matrix rank-nullity then gives $\dim E_{\lambda}(A) =n-\operatorname{rank}(A-\lambda I_n) =n-\operatorname{rank}(A^t-\lambda I_n) =\dim E_{\lambda}(A^t)$.
\begin{enumerate}[label=\roman*.]
\setlength{\itemsep}{0pt}
\item[$(\Rightarrow)$] If $\lambda\in\operatorname{spec}(A)$, then $E_{\lambda}(A)\neq\{\mathbf{0}\}$, so $E_{\lambda}(A^t)\neq\{\mathbf{0}\}$ and $\lambda\in\operatorname{spec}(A^t)$.
\item[$(\Leftarrow)$] The same argument with $A$ and $A^t$ interchanged gives the converse.
\end{enumerate}
\hlt{Structural Commentary}: Transposition preserves rank, so shifted kernels have the same dimension.\\
\hlt{Key Insight}: Spectrum and geometric multiplicity are kernel facts for $A-\lambda I_n$.
\end{proof}
\end{proposition}

\begin{corollary} \label{cor:matrix-invertible-iff-zero-not-spectrum} \hlt{Matrix Invertibility and the Spectrum}\\
Let $A\in\operatorname{Mat}(n,n,\mathbb{F})$. Then $A$ is invertible if and only if $0\notin\operatorname{spec}(A)$.
\begin{proof}
\begin{enumerate}[label=\roman*.]
\setlength{\itemsep}{0pt}
\item[$(\Rightarrow)$] If $A$ is invertible, then $\ker(A)=\{\mathbf{0}\}$, so $0$ is not an eigenvalue of $A$.
\item[$(\Leftarrow)$] If $0\notin\operatorname{spec}(A)$, then $\ker(A)=\{\mathbf{0}\}$. Corollary \ref{cor:matrix-rank-nullity} gives full rank, hence $A$ is invertible.
\end{enumerate}
\hlt{Structural Commentary}: Matrix invertibility is the coordinate form of trivial kernel.\\
\hlt{Key Insight}: A zero eigenvalue is exactly a nonzero solution of $A\mathbf{x}=\mathbf{0}$.
\end{proof}
\end{corollary}

\begin{proposition} \label{prop:trace-determinant-from-eigenvalues} \hlt{Trace and Determinant from Eigenvalues}\\
Let $\lambda_1,\ldots,\lambda_n$ list the eigenvalues of $A\in\operatorname{Mat}(n,n,\mathbb{C})$ according to $m_A$. Then $\operatorname{tr}A=\sum\limits_{j=1}^{n}\lambda_j$ and $\det A=\prod\limits_{j=1}^{n}\lambda_j$.
\begin{proof}
By Fundamental Theorem of Algebra, $p_A(z)=\prod\limits_{j=1}^{n}(z-\lambda_j) =z^n-\left(\sum\limits_{j=1}^{n}\lambda_j\right)z^{n-1}+\cdots+(-1)^n\prod\limits_{j=1}^{n}\lambda_j$.\\
Theorem \ref{thm:matrix-characteristic-polynomial-basic-coefficients} gives $p_A(z)=z^n-(\operatorname{tr}A)z^{n-1}+\cdots+(-1)^n\det A$. Uniqueness of coefficients gives the identities.\\
\hlt{Structural Commentary}: Coefficients of $p_A$ encode symmetric functions of the eigenvalues.\\
\hlt{Key Insight}: Trace is the first root sum and determinant is the full root product.
\end{proof}
\end{proposition}

\begin{proposition} \label{prop:block-upper-triangular-characteristic-polynomial} \hlt{Block Upper-Triangular Characteristic Polynomial}\\
Let $A=
\begin{bmatrix}
B & C\\
0 & D
\end{bmatrix}$, $B$ and $D$ are square. Then $p_A(z)=p_B(z)p_D(z)$, so $m_A(\lambda)=m_B(\lambda)+m_D(\lambda) \ \forall\lambda\in\mathbb{F}$.
\begin{proof}
The shifted matrix is block upper triangular:
\begin{align}
zI-A=
\begin{bmatrix}
zI-B & -C\\
0 & zI-D
\end{bmatrix}.
\nonumber
\end{align}
Theorem \ref{thm:block-upper-triangular-determinant} gives $p_A(z)=\det(zI-B)\det(zI-D)=p_B(z)p_D(z)$.\\
Roots of a product are the roots of the factors, with multiplicity added.\\
\hlt{Structural Commentary}: The off-diagonal block couples vectors but not determinant factors.\\
\hlt{Key Insight}: In block upper-triangular form, spectral data lives on the diagonal blocks.
\end{proof}
\end{proposition}

\begin{theorem} \label{thm:similarity-preserves-characteristic-polynomial} \hlt{Similarity Preserves Characteristic Data}\\
Let $A,B\in\operatorname{Mat}(n,n,\mathbb{F})$ and suppose $B=SAS^{-1}$ for some invertible $S$. Then $p_B=p_A$. In particular, for each $\lambda\in\operatorname{spec}(A)$, $m_B(\lambda)=m_A(\lambda)$ and $\dim E_{\lambda}(B)=\dim E_{\lambda}(A)$.
\begin{proof} $p_B(z) =\det(zI_n-SAS^{-1}) =\det(S(zI_n-A)S^{-1}) =\det(S)\det(zI_n-A)\det(S)^{-1} =p_A(z) \Rightarrow$\\
$m_B(\lambda)=m_A(\lambda)$. Also $B-\lambda I_n=S(A-\lambda I_n)S^{-1}$, multiply by $S$ gives isomorphism $\ker(A-\lambda I_n)\to\ker(B-\lambda I_n)$.\\
\hlt{Structural Commentary}: Similarity changes coordinates, not the shifted determinant or shifted kernel dimension.\\
\hlt{Key Insight}: Conjugation inserts $S$ and $S^{-1}$, which cancel in determinants and identify kernels.
\end{proof}
\end{theorem}

\begin{theorem} \label{thm:geometric-multiplicity-bounded-by-algebraic-multiplicity} \hlt{Geometric Multiplicity Is at Most Algebraic Multiplicity}\\
Let $A\in\operatorname{Mat}(n,n,\mathbb{F})$ and let $\lambda\in\operatorname{spec}(A)$. Then $\dim E_{\lambda}(A)\leq m_A(\lambda)$.
\begin{proof}
Let $k=\dim E_{\lambda}(A)$ and choose a basis $\mathbf{x}_1,\ldots,\mathbf{x}_k$ of $E_{\lambda}(A)$. Extend it to a basis of $\mathbb{F}^n$, and let $S$ be the matrix whose columns are this basis. In the new basis,
\begin{align}
S^{-1}AS=
\begin{bmatrix}
\lambda I_k & *\\
0 & C
\end{bmatrix}.
\nonumber
\end{align}
By Theorem \ref{thm:similarity-preserves-characteristic-polynomial} and Proposition \ref{prop:block-upper-triangular-characteristic-polynomial}, $p_A(z)=p_{S^{-1}AS}(z)=(z-\lambda)^kp_C(z)$. Thus $m_A(\lambda)\geq k=\dim E_{\lambda}(A)$.\\
\hlt{Structural Commentary}: An eigenspace-first basis exposes a $\lambda I_k$ block.\\
\hlt{Key Insight}: Each independent eigenvector forces one visible factor $z-\lambda$ in the characteristic polynomial.
\end{proof}
\end{theorem}

\begin{theorem} \label{thm:diagonalizable-iff-algebraic-geometric-multiplicities-match} \hlt{Diagonalizability and Multiplicity}\\
Let $A\in\operatorname{Mat}(n,n,\mathbb{C})$. Then $A$ is diagonalizable if and only if $\dim E_{\lambda}(A)=m_A(\lambda)$ for every $\lambda\in\operatorname{spec}(A)$.
\begin{proof}
\begin{enumerate}[label=\roman*.]
\setlength{\itemsep}{0pt}
\item[$(\Rightarrow)$] If $A$ is diagonalizable, then $A$ is similar to a diagonal matrix $D$. Theorem \ref{thm:similarity-preserves-characteristic-polynomial} gives $m_A(\lambda)=m_D(\lambda)$ and $\dim E_{\lambda}(A)=\dim E_{\lambda}(D)$. For diagonal $D$, both numbers count diagonal entries equal to $\lambda$.
\item[$(\Leftarrow)$] Assume $\dim E_{\lambda}(A)=m_A(\lambda)$ for every eigenvalue.\\
If $\mu_1,\ldots,\mu_d$ are distinct eigenvalues, then $\sum\limits_{i=1}^{d}m_A(\mu_i)=n$. Hence $\dim E_{\mu_1}(A)+\cdots+\dim E_{\mu_d}(A)=n$.\\
Theorem \ref{thm:conditions-equivalent-to-diagonalizability} gives an eigenbasis, so $A$ is diagonalizable.
\end{enumerate}
\hlt{Structural Commentary}: Diagonalizability means root counts match eigenspace dimensions.\\
\hlt{Key Insight}: No missing eigendirections means an eigenbasis.
\end{proof}
\end{theorem}

\begin{theorem} \label{thm:polynomial-functional-calculus-diagonalizable-matrix} \hlt{Polynomial Functional Calculus for Diagonalizable Matrices}\\
Let $A\in\operatorname{Mat}(n,n,\mathbb{F})$ be diagonalizable and write $A=S\Lambda S^{-1}$, where $\Lambda=\operatorname{diag}(\lambda_1,\ldots,\lambda_n)$.\\
If $p\in\mathcal{P}(\mathbb{F})$, then $p(A)=S\operatorname{diag}(p(\lambda_1),\ldots,p(\lambda_n))S^{-1}$.
\begin{proof}
For every $j\geq0$, $A^j=S\Lambda^jS^{-1}$.\\
If $p(z)=\sum_j c_jz^j$, then $p(A)=\sum\limits_j c_jA^j =S\left(\sum\limits_j c_j\Lambda^j\right)S^{-1} =S p(\Lambda) S^{-1}$.\\
Since $\Lambda$ is diagonal, $p(\Lambda)=\operatorname{diag}(p(\lambda_1),\ldots,p(\lambda_n))$.\\
\hlt{Structural Commentary}: Diagonalization reduces $p(A)$ to scalar evaluation on eigenvalues.\\
\hlt{Key Insight}: Compute functions of $A$ by applying the function to the eigenvalues.
\end{proof}
\end{theorem}

\begin{definition} \label{def:matrix-commutant} \hlt{Commutant}\\
For $\emptyset\neq\mathcal{F}\subseteq\operatorname{Mat}(n,n,\mathbb{F})$, the \hlt{commutant} of $\mathcal{F}$ is $\mathcal{F}'\defeq\{X\in\operatorname{Mat}(n,n,\mathbb{F}): XA=AX\ \forall A\in\mathcal{F}\}$.\\
For a single matrix $A$, write $\{A\}'$.
\end{definition}

\begin{theorem} \label{thm:commutant-of-diagonalizable-matrix} \hlt{Commutant of a Diagonalizable Matrix}\\
Let $A\in\operatorname{Mat}(n,n,\mathbb{C})$ be diagonalizable. Suppose its distinct eigenvalues are $\mu_1,\ldots,\mu_d$, and set $n_i=m_A(\mu_i)$.\\
Write $A=S(\mu_1I_{n_1}\oplus\cdots\oplus\mu_dI_{n_d})S^{-1}$. Then $\{A\}'=\{S(X_1\oplus\cdots\oplus X_d)S^{-1}:X_i\in\operatorname{Mat}(n_i,n_i,\mathbb{C})\}$.\\
In particular, $\dim\{A\}'=n_1^2+\cdots+n_d^2$.
\begin{proof}
Let $\Lambda=\mu_1I_{n_1}\oplus\cdots\oplus\mu_dI_{n_d}$. Matrix $B$ commutes with $A$ if and only if $S^{-1}BS$ commutes with $\Lambda$.\\
By Lemma \ref{lem:block-diagonal-commutant-distinct-scalars}, this happens exactly when $S^{-1}BS=X_1\oplus\cdots\oplus X_d$ with $X_i\in\operatorname{Mat}(n_i,n_i,\mathbb{C})$.\\
Dimension count follows because conjugation by $S$ is a vector-space isomorphism and $\dim\operatorname{Mat}(n_i,n_i,\mathbb{C})=n_i^2$.\\
\hlt{Structural Commentary}: Diagonalizing $A$ reduces commutation to respecting equal-eigenvalue blocks.\\
\hlt{Key Insight}: Distinct eigenvalue blocks cannot mix under a commuting matrix.
\end{proof}
\end{theorem}

\begin{corollary} \label{cor:commutant-distinct-eigenvalues-polynomials} \hlt{Simple Spectrum Commutants Are Polynomials}\\
If $A\in\operatorname{Mat}(n,n,\mathbb{C})$ has $n$ distinct eigenvalues, then $\{A\}'=\{p(A):p\in\mathcal{P}(\mathbb{C})\}$, $\dim\{A\}'=n$.
\begin{proof}
Write $A=S\operatorname{diag}(\lambda_1,\ldots,\lambda_n)S^{-1}$ with distinct $\lambda_i$. Theorem \ref{thm:commutant-of-diagonalizable-matrix} says every $B\in\{A\}'$ has the form $S\operatorname{diag}(\xi_1,\ldots,\xi_n)S^{-1}$. Choose the Lagrange interpolation polynomial $p(z)=\sum\limits_{i=1}^{n}\xi_i\prod\limits_{j\neq i}\frac{z-\lambda_j}{\lambda_i-\lambda_j}$.\\
Then $p(\lambda_i)=\xi_i$ for every $i$, so Theorem \ref{thm:polynomial-functional-calculus-diagonalizable-matrix} gives $B=p(A)$.\\
The dimension statement is the case $n_1=\cdots=n_n=1$ of Theorem \ref{thm:commutant-of-diagonalizable-matrix}.\\
\hlt{Structural Commentary}: Simple spectrum leaves one scalar choice per eigendirection.\\
\hlt{Key Insight}: Interpolation converts arbitrary diagonal data into a polynomial in $A$.
\end{proof}
\end{corollary}
